\chapter{Methods}


% \item The vector elements in short: 
%     \begin{enumerate}
%     \item The spectral slope ($\beta$) is a first-order summary of basal character. High $\beta$ indicates smooth, soft beds (sedimentary basins, active ice streams), while low $\beta$ indicates hard, jagged crystalline shields (e.g., Jordan et al. (2017)\cite{Jordan_2017}; Bingham et al (2009)\cite{Bingham_2009}). 
%     \item Roughness anisotropy ($\Delta\beta$): Topographic roughness varies relative to ice flow. Smoothness along-flow implies glacial erosion and streamlining, while cross-flow roughness often indicates geological structure like rifts (e.g., Hubbard et al (2000)\cite{Hubbard_2000}; Rippin et al. 2014 \cite{Rippin_2014}).



%     \item
    
%     \item

%     \item DESCRIPTORS

%         \item Roughness amplitude ($A_{1\text{km}}$), is necessary to separate features like flat crystalline plains from mountainous rift flanks (e.g., Cooper et al (2019)\cite{Cooper_2019}).

%             \item 

%             \item 

%             \item 

%             \item 

%             \item 

%             \item 

%             \item 
            

%     \end{enumerate}